\chapter{Component Description and Contribution to Objectives\label{chap:components}}

This chapter provides a detailed description of each system component and explains how it contributes to the project objectives defined in Chapter~\ref{chap:introduction}.

\section{Custom Action Interface: CatchTarget}

\subsection{Description}

The \texttt{CatchTarget} action is the central communication contract between \texttt{master\_manager} and \texttt{catch\_executor}. It is defined in \texttt{CatchTarget.action} as follows:

\begin{lstlisting}[language={},title={CatchTarget.action}]
string target_name
---
bool success
string caught_name
---
float32 distance_remaining
\end{lstlisting}

The three parts of the action definition serve distinct purposes:

\begin{itemize}
  \item \textbf{Goal}: \texttt{target\_name} specifies which turtle to pursue, identified by its name (e.g., \texttt{turtle2}).
  \item \textbf{Result}: \texttt{success} indicates whether the catch succeeded; \texttt{caught\_name} confirms the identity of the caught turtle, providing an unambiguous result even in edge cases.
  \item \textbf{Feedback}: \texttt{distance\_remaining} provides real-time distance to the target during execution, enabling the master manager to monitor progress and make preemption decisions.
\end{itemize}

\subsection{Contribution to Objectives}

This action interface directly supports \textbf{Objective~O2} (Action-Based Task Execution) by providing a structured, asynchronous mechanism for long-running pursuit tasks. The feedback channel enables the master manager to monitor progress and make preemption decisions, while the result channel ensures reliable outcome reporting. The use of a custom action interface---rather than a simple service---demonstrates proper use of the ROS2 Action paradigm for tasks that are long-running, provide feedback, and can be cancelled.

\section{Utility Module: utils.py}

\subsection{Description}

The \texttt{utils.py} module provides shared mathematical helper functions used across all nodes. These functions encapsulate common geometric computations that are fundamental to the control and decision-making logic:

\begin{itemize}
  \item \texttt{normalize\_angle(theta)}: Wraps an angle to the $[-\pi, \pi]$ range. This is essential for correct heading error computation, as raw angle differences can exceed $2\pi$ and lead to incorrect turning directions.
  \item \texttt{distance(x1, y1, x2, y2)}: Computes the Euclidean distance between two points:
  \[
    d = \sqrt{(x_2 - x_1)^2 + (y_2 - y_1)^2}
  \]
  \item \texttt{angle\_to(x\_from, y\_from, x\_to, y\_to)}: Computes the bearing angle from one point to another:
  \[
    \theta = \text{atan2}(y_{to} - y_{from},\ x_{to} - x_{from})
  \]
  \item \texttt{clamp(value, lo, hi)}: Limits a value within a specified range, used to enforce velocity limits.
  \item \texttt{sign(value)}: Returns $+1$, $-1$, or $0$ based on the sign of the input, used to determine rotation direction.
\end{itemize}

\subsection{Contribution to Objectives}

This module supports \textbf{all objectives} by ensuring consistent and correct geometric computations. Without proper angle normalization, the control logic would produce incorrect turning commands; without distance calculations, target selection and catch detection would be impossible. Centralising these functions in a shared module also eliminates code duplication and reduces the risk of inconsistent implementations across nodes.

\section{Turtle Registry: turtle\_registry.py}

\subsection{Description}

The \texttt{TurtleRegistry} class provides a centralised data structure for tracking all turtles known to the system. It stores the following information for each turtle via the \texttt{TurtleEntry} dataclass:

\begin{itemize}
  \item \texttt{name}: Turtle identifier (e.g., \texttt{turtle2}).
  \item \texttt{x}, \texttt{y}: Current position coordinates.
  \item \texttt{caught}: Boolean flag indicating whether the turtle has been caught.
  \item \texttt{has\_pose}: Boolean flag indicating whether a pose message has been received.
\end{itemize}

Key methods include:

\begin{itemize}
  \item \texttt{add(name)}: Registers a new turtle in the registry.
  \item \texttt{update\_pose(name, x, y)}: Updates a turtle's position from pose messages.
  \item \texttt{mark\_caught(name)}: Marks a turtle as caught after successful pursuit.
  \item \texttt{uncaught\_targets(exclude)}: Returns all turtles that have poses and are not yet caught, optionally excluding specific names.
  \item \texttt{nearest\_to(x, y, exclude)}: Finds the nearest uncaught turtle to a given position, using Euclidean distance as the metric.
\end{itemize}

\subsection{Contribution to Objectives}

The registry directly supports \textbf{Objective~O1} (Autonomous Target Discovery and Pursuit) by maintaining a live state table of all turtles. It enables the master manager to efficiently query for uncaught targets and select the nearest one, which is the core decision-making primitive of the system. The separation of the registry into its own module follows the single-responsibility principle and makes the state management logic testable independently of the ROS2 node logic.

\section{Spawn Manager: spawn\_manager.py}

\subsection{Description}

The \texttt{spawn\_manager} node is responsible for periodically generating new target turtles in the simulation. It operates as a service client that calls the \texttt{/spawn} service provided by \texttt{turtlesim\_node}.

\textbf{Key Features:}

\begin{enumerate}
  \item \textbf{Periodic Spawning}: A timer fires every \texttt{spawn\_period} seconds (default: 3.0\,s) to create a new turtle at a random position within configurable bounds.

  \item \textbf{Automatic Index Management}: On startup, the node scans existing \texttt{/turtleN/pose} topics using a regular expression pattern to determine the next available index. This prevents name collisions after node restarts, as the node will always start numbering from $\max(\text{existing indices}) + 1$.

  \item \textbf{Duplicate Name Handling}: If \texttt{turtlesim} rejects a spawn request (returns an empty name string), the node automatically increments the index and retries on the next timer tick. This handles edge cases where turtles may have been spawned externally.

  \item \textbf{Failure Resilience}: After a configurable number of consecutive failures (\texttt{max\_consecutive\_failures}, default: 5), the node re-scans existing topics and resynchronises its index counter, recovering from situations where the index has drifted out of sync with reality.
\end{enumerate}

\subsection{Configurable Parameters}

\begin{table}[htbp]
\centering
\caption{Spawn Manager Parameters}
\begin{tabular}{llc}
\toprule
\textbf{Parameter} & \textbf{Description} & \textbf{Default} \\
\midrule
\texttt{spawn\_period} & Time between spawns (s) & 3.0 \\
\texttt{x\_min}, \texttt{x\_max} & X-coordinate bounds & 1.0, 10.0 \\
\texttt{y\_min}, \texttt{y\_max} & Y-coordinate bounds & 1.0, 10.0 \\
\texttt{start\_index} & First turtle number & 2 \\
\texttt{max\_consecutive\_failures} & Failure threshold for rescan & 5 \\
\bottomrule
\end{tabular}
\end{table}

\subsection{Contribution to Objectives}

This node directly implements \textbf{Objective~O3} (Dynamic Environment Generation). By continuously introducing new targets at random positions, it creates a dynamic and unpredictable environment that forces the rest of the system to continuously adapt its pursuit strategy. The robustness features (index scanning, duplicate handling, failure resilience) ensure that the system can run indefinitely without manual intervention, contributing to \textbf{Objective~O5} (Robust System Design).

\section{Master Manager: master\_manager.py}

\subsection{Description}

The \texttt{master\_manager} node is the central decision-making component---the ``brain''---of the system. It discovers turtles, selects pursuit targets, dispatches catch goals via the Action client, and manages the caught turtle chain.

\textbf{Key Features:}

\begin{enumerate}
  \item \textbf{Auto-Discovery}: Periodically scans the ROS2 topic list for \texttt{/turtleN/pose} topics and automatically subscribes to new turtles' poses. This eliminates the need for explicit hand-off between \texttt{spawn\_manager} and \texttt{master\_manager}, achieving loose coupling between the two nodes.

  \item \textbf{Nearest-Target Selection}: Every \texttt{decision\_period} (default: 0.5\,s), the node evaluates all uncaught turtles and selects the one closest to the master turtle using Euclidean distance.

  \item \textbf{Action Client}: Sends \texttt{CatchTarget} goals to \texttt{catch\_executor} and handles results, including success, failure, and cancellation.

  \item \textbf{Preemption with Hysteresis}: While a catch goal is in flight, the node continuously re-evaluates whether a closer target has appeared. A preemption occurs only when:
  \[
    d_{\text{new}} + m_{\text{preempt}} < d_{\text{current}}
  \]
  where $m_{\text{preempt}}$ is the \texttt{preempt\_margin} parameter (default: 1.0\,m). This hysteresis prevents rapid flip-flopping between two nearly equidistant targets, which would cause erratic behaviour.

  \item \textbf{Failure Cooldown}: When a catch goal fails (not due to preemption), the target is placed on a cooldown list for \texttt{failure\_cooldown\_sec} seconds (default: 5.0\,s), preventing the system from repeatedly attempting to catch a problematic target.

  \item \textbf{Chain Management}: Upon successful catch, the turtle is appended to the internal chain list and the updated chain is published on the \texttt{/caught\_chain} topic as a JSON message containing the leader name and the ordered chain of caught turtles.
\end{enumerate}

\subsection{Configurable Parameters}

\begin{table}[htbp]
\centering
\caption{Master Manager Parameters}
\begin{tabular}{llc}
\toprule
\textbf{Parameter} & \textbf{Description} & \textbf{Default} \\
\midrule
\texttt{master\_name} & Name of the master turtle & \texttt{turtle1} \\
\texttt{discover\_period} & Topic scan interval (s) & 1.0 \\
\texttt{decision\_period} & Decision tick interval (s) & 0.5 \\
\texttt{failure\_cooldown\_sec} & Cooldown for failed targets (s) & 5.0 \\
\texttt{action\_server\_wait\_sec} & Action server wait timeout (s) & 2.0 \\
\texttt{preempt\_margin} & Preemption hysteresis margin (m) & 1.0 \\
\bottomrule
\end{tabular}
\end{table}

\subsection{Contribution to Objectives}

The master manager is the most critical component for achieving \textbf{Objective~O1} (Autonomous Target Discovery and Pursuit) and \textbf{Objective~O2} (Action-Based Task Execution). It orchestrates the entire pursuit workflow: discovering targets, making intelligent decisions about which target to pursue, managing the action lifecycle, and coordinating the formation chain. The preemption mechanism ensures that the system always pursues the most efficient target, while the cooldown mechanism prevents pathological retry loops, both contributing to \textbf{Objective~O5} (Robust System Design).

\section{Catch Executor: catch\_executor.py}

\subsection{Description}

The \texttt{catch\_executor} node implements the Action Server for the \texttt{catch\_target} action. It physically controls the master turtle's velocity to pursue and catch a designated target.

\textbf{Key Features:}

\begin{enumerate}
  \item \textbf{Single-Goal Serialization}: While one goal is executing, any new goal request is \texttt{REJECTED}. This ensures that two control loops never simultaneously publish conflicting velocity commands to \texttt{/turtle1/cmd\_vel}. Preemption is handled by the master manager canceling the current goal, not by queuing multiple goals.

  \item \textbf{Control Law}: At each control tick (default: 20\,Hz), the node computes the distance and heading error to the target and applies a turn-then-drive strategy (detailed in Chapter~\ref{chap:algorithms}).

  \item \textbf{Catch Detection}: When the distance falls below \texttt{catch\_distance} (default: 0.5\,m), the goal succeeds and the result is returned.

  \item \textbf{Timeout Handling}: Two timeout mechanisms protect against stuck goals:
  \begin{itemize}
    \item \texttt{goal\_timeout\_sec} (default: 30.0\,s): Total execution time limit.
    \item \texttt{no\_pose\_timeout\_sec} (default: 5.0\,s): Maximum wait for pose data before aborting.
  \end{itemize}

  \item \textbf{Cancellation Support}: Goals can be cancelled at any time, and the node immediately stops the turtle by publishing a zero-velocity command.
\end{enumerate}

\subsection{Configurable Parameters}

\begin{table}[htbp]
\centering
\caption{Catch Executor Parameters}
\begin{tabular}{llc}
\toprule
\textbf{Parameter} & \textbf{Description} & \textbf{Default} \\
\midrule
\texttt{master\_name} & Name of the master turtle & \texttt{turtle1} \\
\texttt{linear\_speed} & Forward/backward speed (m/s) & 1.5 \\
\texttt{angular\_speed} & Turning speed (rad/s) & 3.0 \\
\texttt{max\_angular\_speed} & Maximum angular speed (rad/s) & 4.0 \\
\texttt{catch\_distance} & Catch threshold (m) & 0.5 \\
\texttt{angle\_tolerance} & Heading tolerance (rad) & 0.1 \\
\texttt{control\_rate\_hz} & Control loop frequency (Hz) & 20.0 \\
\texttt{goal\_timeout\_sec} & Total goal timeout (s) & 30.0 \\
\texttt{no\_pose\_timeout\_sec} & No-pose timeout (s) & 5.0 \\
\texttt{allow\_reverse} & Enable backward driving & False \\
\bottomrule
\end{tabular}
\end{table}

\subsection{Contribution to Objectives}

The catch executor directly implements \textbf{Objective~O2} (Action-Based Task Execution) by providing a robust Action Server that handles the physical pursuit task. The single-goal serialization ensures system stability by preventing conflicting velocity commands, while the timeout mechanisms and cancellation support contribute to \textbf{Objective~O5} (Robust System Design). The optional reverse driving feature (\texttt{allow\_reverse}) demonstrates extensibility in control strategy design.

\section{Follower Manager: follower\_manager.py}

\subsection{Description}

The \texttt{follower\_manager} node implements the chain formation behaviour. After a turtle is caught, it is added to a following chain where each caught turtle follows its predecessor.

\textbf{Key Features:}

\begin{enumerate}
  \item \textbf{Chain Subscription}: Listens to the \texttt{/caught\_chain} topic for JSON messages containing the leader name and the ordered chain of caught turtles.

  \item \textbf{Dynamic Subscription Management}: Automatically creates pose subscribers and velocity publishers for each new follower as it joins the chain.

  \item \textbf{Proportional Following Control}: At each control tick (default: 20\,Hz), for each follower in the chain, the node computes the distance and heading error to its leader and applies proportional control (detailed in Chapter~\ref{chap:algorithms}).

  \item \textbf{Chain Rule}: The first caught turtle follows \texttt{turtle1}, the second follows the first, the third follows the second, and so on, forming a snake-like chain.
\end{enumerate}

\subsection{Configurable Parameters}

\begin{table}[htbp]
\centering
\caption{Follower Manager Parameters}
\begin{tabular}{llc}
\toprule
\textbf{Parameter} & \textbf{Description} & \textbf{Default} \\
\midrule
\texttt{leader\_name} & Chain leader name & \texttt{turtle1} \\
\texttt{follow\_distance} & Desired following distance (m) & 0.8 \\
\texttt{linear\_speed} & Maximum linear speed (m/s) & 1.2 \\
\texttt{angular\_speed} & Maximum angular speed (rad/s) & 2.5 \\
\texttt{max\_angular\_speed} & Angular speed cap (rad/s) & 3.5 \\
\texttt{angle\_tolerance} & Heading tolerance (rad) & 0.15 \\
\texttt{control\_rate\_hz} & Control loop frequency (Hz) & 20.0 \\
\texttt{linear\_kp} & Proportional gain for linear speed & 1.0 \\
\texttt{angular\_kp} & Proportional gain for angular speed & 4.0 \\
\bottomrule
\end{tabular}
\end{table}

\subsection{Contribution to Objectives}

The follower manager directly implements \textbf{Objective~O4} (Chain Formation Following). By maintaining a dynamic chain of caught turtles and applying proportional control for each follower, it creates a visually compelling snake-like formation. The proportional control with distance threshold prevents oscillation and collision, contributing to \textbf{Objective~O5} (Robust System Design).

\section{Launch File: catch\_turtle.launch.py}

\subsection{Description}

The launch file starts all five nodes with a single command and loads the shared parameter file from \texttt{config/params.yaml}. Each node is configured with the appropriate package, executable, and parameter file path.

\subsection{Contribution to Objectives}

The launch file provides a single-command startup mechanism that ensures all nodes are launched in the correct configuration with shared parameters. This simplifies deployment and testing, and ensures that the system starts in a consistent state every time.
